\subsection{Introduction to Vector}
In Grade 12 Calculus and Vector, we will informally dicuss "Vector" and its basic operatoins and applications.
\begin{theorem}
    [Vector] A vector is a quantitiy that has both magnitude and direction. 
\end{theorem} 
\begin{example}
    A statement \textbf{A car travels $120km/h$ heading northeast} can be represented by a vector. 
    \begin{equation*}
        \vec{v} = 120\tfrac{km}{h}\left [ \text{North $45^{\circ}$ of East}\right ]
    \end{equation*}
\end{example}

\subsubsection{Standard Convention}
There are three ways to measure the direction of a vector.
\begin{itemize}
    \item Angles measured from Standard Position
    \begin{itemize}
        \item starts from a horizontal position (positive arm of x-axis)
        \item single measured by rotating counterclockwise
    \end{itemize}
    \item Angles established using a true bearing
    \begin{itemize}
        \item always start from North, then rotate clockwise. 
        \item express all angles using 3 digits
    \end{itemize}
    \item Angles established using a quadrant bearing
    \begin{itemize}
        \item start from either a North or South position
        \item rotate East or West by an amount in between $0^{\circ}$ and $360^{\circ}$
    \end{itemize}
\end{itemize}

\subsubsection{Common terminology}
Below are common terminologies in vector. 
\begin{enumerate}
    \item Parallel vectors 
    \begin{itemize}
        \item vector point in the same or opposite direction 
        \item no relationship between their sizes or magnitudes
    \end{itemize}
    \item Equivalent vectors 
    \begin{itemize}
        \item vectors that are identical (size and direction)
        \item these do not have to be at the same "location"
    \end{itemize}
    \item Opposite vectors 
    \begin{itemize}
        \item these vectors are equal in magnitude 
        \item directinos are opposite ($180^{\circ}$ away from the other)
    \end{itemize}
\end{enumerate}

\subsubsection{Addition Vectors}
\begin{theorem}
    [Commutative property of Vectors]
    Given vectors $\vec{u}$ and $\vec{v}$ $$\vec{v} + \vec{u} = \vec{u} + \vec{v}$$
\end{theorem}

Geometrically, the addition of vectors requires arranging the vectors from tip of one vector to the tail of the other. 
This method is referred to as the TriangleMethod.

For TriangleMethod, given $\vec{a}$ and $\vec{b}$, define $\vec{c} = \vec{a} + \vec{b}$, following statements are always true:
\begin{itemize}
    \item $\vec{c} = \vec{a} + \vec{b}$
\end{itemize}

The following statements are not always true:
\begin{itemize}
    \item $\left|\vec{c}\right| = \left|\vec{a}\right| + \left|\vec{b}\right|$
    \item $(\left|\vec{c}\right|)^2 = (\left|\vec{a}\right|)^2 + (\left|\vec{b}\right|)^2$
\end{itemize}

Instead of convert vertors into matrixes, do matrixes addition, and convert the resultant matrix back to vector. 
There are also other ways to add vectors. \\

Given vectors $\vec{v_1}$, $\vec{v_2}$ and resultant vector $\vec{v_r}$,
if $\vec{v_1}$ and $\vec{v_2}$ are parallel vector to each other:
\[
    \left|\vec{v_r}\right| = \left|\vec{v_1}\right| + \left|\vec{v_2}\right|
\]
If the included angle between $\vec{v_1}$ and $\vec{v_2}$ is $90^{\circ}$, we can use Pythagorean theorem:
\[
    \left|\vec{v_r}\right| = \sqrt{\left|v_1\right|^2 - \left|v_2\right|^2}
\]
Otherwise, we can use COSINE Law:
\[
    \left|\vec{v_r}\right| = \sqrt{\left|v_1\right|^2 + \left|v_2\right|^2 - 2\left|v_1\right|\left|v_2\right|\cos \theta}
\]

\subsubsection{Subtraction of vectors}
When subtracting vectors arrange them tail-to-tail. The resultant vector should be oriented so that the tail (starting point) goes from the head of the second 
vector to the head of the first vector, which will be the end point, as indicated by the "difference of vectors" expression. 

\newpage